\section{Introduzione}

KompozeR è una web application che estende le funzionalità dell'e-commerce di Kompo (www.soluzionekompo.com), azienda specializzata in mobili componibili personalizzabili. 
L'applicazione permette all'utente di progettare la propria scaffalatura tramite un configuratore visuale interattivo, ottenendo un'anteprima del risultato finale e popolando automaticamente il carrello con i componenti selezionati.
L'utente oltre alla creazione del proprio scaffale personalizzato potrà accedere a funzionalità avanzate di assistenza e monitoraggio in tempo reale. 
Un chatbot assistente in tempo reale guida l'utente durante la configurazione, rispondendo a domande su dimensioni, compatibilità e prezzi. 
Il sistema notifica inoltre l'utente in tempo reale in caso di variazioni di disponibilità o prezzo dei componenti inseriti nella propria configurazione.
Lato amministratore, KompozeR fornisce strumenti per la gestione del catalogo prodotti, un sistema di monitoraggio degli ordini (che traccia l'arrivo e la spedizione dei componenti) e una pagina di reportistica per vedere i trend di vendita.

\subsection{Contesto}
Kompo è un'azienda che si occupa di mobili componibili personalizzabili, offrendo soluzioni su misura per ogni esigenza abitativa.
Il catalogo si divide in 3 categorie di prodotti: Tondo, Quuadro e Qube; non compatibili tra loro, ma con caratteristiche e funzionalità simili 
(Questo dettaglio è importante per definire in seguito la logica di componimento dei singoli componenti: che è la medesima per ogni tipologia di prodotto).
Il processo di acquisto attuale prevede che l'utente selezioni i singoli componenti desiderati, avendo già un'idea chiara del risultato finale, e li aggiunga al carrello.
Questo processo può risultare complesso e poco intuitivo per gli utenti che non conoscono il prodotto e si trovano costretti a chiedere assestenza direttamente al personale di vendita, con conseguente aumento dei tempi di acquisto e potenziale insoddisfazione del cliente.
KompozeR nasce con l'obiettivo di semplificare e migliorare l'esperienza di acquisto, offrendo un configuratore visuale il più semplice e intuitivo possibile, garantento a tutti la possibilità di acquisto.
